% Source PDF: source/普通高考/2013/2013辽宁理.pdf, page 1, question 8.
% Grid evidence: tmp/review_2013_liaoning_science/grid/page-1-grid100.jpg and
%   q8_original_clean2.png (no-grid crop from the source page).
% Elements: rounded start/end terminals; input parallelogram; initialization
% rectangle S=0,i=2; decision diamond i<=n; update rectangles
% S=S+1/(i^2-1), i=i+2; output parallelogram. Complete node -> directed-edge list:
% start -> input -> init -> test; test --是--> update -> increment -> left loop
% -> main merge before test; test --否--> output -> finish. The return arrow,
% main downward arrow, and output branch are independent; all connectors solid,
% no other edges. Node text is centered with local zero paragraph indentation.

\begin{tikzpicture}[
  >=Stealth,
  x=1cm,y=0.88cm,
  line width=0.45pt,line cap=round,line join=round,
  font=\normalsize,
  flowtext/.style={anchor=center,align=center,inner sep=0pt,
    execute at begin node={\setlength{\parindent}{0pt}}},
  every node/.style={flowtext},
  flowbox/.style={flowtext,draw,minimum height=.68cm,inner xsep=3pt,inner ysep=2pt},
  terminal/.style={flowbox,rounded corners=.16cm,minimum width=1.35cm},
  io/.style={flowbox,shape=trapezium,trapezium left angle=72,trapezium right angle=72,
    minimum height=.76cm},
    process/.style={flowbox},
    decision/.style={flowbox,diamond,aspect=2.8,minimum width=3.75cm,minimum height=1.04cm,inner sep=0pt},
    branchlabel/.style={flowtext,inner sep=0pt},
]
  \path[use as bounding box] (-3.45,0.85) rectangle (4.75,10.15);
  \node[terminal] (start) at (0,9.65) {开始};
  \node[io,minimum width=2.05cm] (input) at (0,8.35) {输入 $n$};
  \node[process,minimum width=3.55cm] (init) at (0,7.00) {$S=0,i=2$};
  \node[decision] (test) at (0,5.48) {$i\le n$};
  \node[process,minimum width=4.05cm,minimum height=1.00cm] (update) at (0,3.90)
    {$S=S+\dfrac{1}{i^2-1}$};
  \node[process,minimum width=2.55cm] (inc) at (0,2.42) {$i=i+2$};
  \node[io,minimum width=2.05cm] (output) at (3.55,4.00) {输出 $S$};
  \node[terminal] (finish) at (3.55,2.42) {结束};

  % Main trunk and terminating false branch.
  \draw[->] (start.south) -- (input.north);
  \draw[->] (input.south) -- (init.north);
  \draw[->] (init.south) -- (test.north);
  \draw[->] (test.east) -- node[branchlabel,above] {否} (3.55,5.48) -- (3.55,4.40) -- (output.north);
  \draw[->] (output.south) -- (finish.north);

  % True update path and return loop.
  \draw[->] (test.south) -- node[branchlabel,right] {是} (update.north);
  \draw[->] (update.south) -- (inc.north);
  \coordinate (loop-bottom) at (0,1.15);
  \coordinate (loop-left) at (-3.15,1.15);
  \coordinate (loop-entry) at (-3.15,6.30);
  \draw (inc.south) -- (loop-bottom) -- (loop-left) -- (loop-entry);
  \draw[->] (loop-entry) -- (0,6.30);
\end{tikzpicture}
